\documentclass[12pt,a4paper]{article}

% ---------- Pacotes ----------
\usepackage[utf8]{inputenc}
\usepackage[T1]{fontenc}
\usepackage[brazilian]{babel}
\usepackage{lmodern}
\usepackage[margin=2.5cm]{geometry}
\usepackage{microtype}
\usepackage{setspace}
\usepackage{xcolor}
\usepackage{titlesec}
\usepackage{enumitem}
\usepackage[most]{tcolorbox}
\usepackage{fancyhdr}
\usepackage[hidelinks]{hyperref}

\onehalfspacing
\setlength{\parindent}{0pt}
\setlength{\parskip}{0.5em}

% ---------- Cores ----------
\definecolor{azul}{HTML}{1F3A5F}
\definecolor{azulclaro}{HTML}{EAF0F8}
\definecolor{cinza}{HTML}{555555}

% ---------- Títulos de seção ----------
\titleformat{\section}
  {\Large\bfseries\color{azul}}
  {\thesection.}{0.6em}{}
\titlespacing*{\section}{0pt}{1.8em}{0.8em}

% ---------- Cabeçalho e rodapé ----------
\pagestyle{fancy}
\fancyhf{}
\renewcommand{\headrulewidth}{0.4pt}
\lhead{\small\color{cinza} Segurança Digital}
\rhead{\small\color{cinza} Atividade --- Bug Bounty e Divulgação de Vulnerabilidades}
\cfoot{\small\color{cinza} \thepage}

% ---------- Perguntas e respostas ----------
% Numeração das perguntas: <seção>.<pergunta>
\newcounter{pergunta}[section]
\renewcommand{\thepergunta}{\thesection.\arabic{pergunta}}

% \pergunta{texto}{resposta}  (resposta pode ficar vazia)
\newcommand{\pergunta}[2]{%
  \refstepcounter{pergunta}%
  \noindent\textbf{\color{azul}\thepergunta}\hspace{0.5em}\textbf{#1}\par
  \noindent\textbf{Resposta:} #2\par
}

% ---------- Documento ----------
\begin{document}

% ----- Cabeçalho do trabalho -----
\begin{center}
  {\LARGE\bfseries\color{azul} Atividade: Bug Bounty, Divulgação\\[2pt] Responsável e Classificação de Vulnerabilidades}\\[1.2em]
  {\large Segurança Digital --- Engenharia de Software}
\end{center}

\vspace{1em}

% ======================================================
\section{O que é Bug Bounty?}

\pergunta{O que significa Bug Bounty?}{}

\pergunta{Por que empresas pagam pessoas para encontrar vulnerabilidades?}{}

\pergunta{O que é uma vulnerabilidade de segurança?}{}

\pergunta{Qual é a diferença entre encontrar uma falha com autorização e explorar uma falha sem autorização?}{}

\pergunta{Cite duas empresas que possuem programas de Bug Bounty.}{}

\pergunta{O que uma pessoa precisa fazer para participar de um programa desses?}{}

% ======================================================
\section{Meta Bug Bounty --- como funciona na prática?}

\pergunta{O que é o programa?}{}

\pergunta{Que tipos de vulnerabilidades podem ser reportadas?}{}

\pergunta{Como uma vulnerabilidade deve ser comunicada à empresa?}{}

\pergunta{O que pode acontecer quando um pesquisador encontra uma falha válida?}{}

\pergunta{Por que a empresa oferece recompensas financeiras?}{}

\pergunta{Encontre um exemplo público de vulnerabilidade reportada à Meta.}{}

% ======================================================
\section{Discord e reconhecimento de pesquisadores}

\pergunta{O Discord possui programa de Bug Bounty?}{}

\pergunta{Como pesquisadores podem reportar vulnerabilidades?}{}

\pergunta{Existem formas de reconhecimento além de dinheiro?}{}

\pergunta{O que são badges ou outras formas de reconhecimento para pesquisadores?}{}

\pergunta{Por que uma empresa pode querer criar uma comunidade de pesquisadores de segurança?}{}

% ======================================================
\section{O que é Responsible Disclosure?}

\pergunta{O que significa Responsible Disclosure?}{}

\pergunta{Por que o pesquisador deve comunicar a empresa antes de divulgar publicamente uma vulnerabilidade?}{}

\pergunta{O que pode acontecer se uma vulnerabilidade for divulgada antes que exista uma correção?}{}

\pergunta{Qual é a diferença entre Responsible Disclosure e simplesmente publicar uma vulnerabilidade na internet?}{}

% ======================================================
\section{CVE: o ``RG'' das vulnerabilidades}

\pergunta{O que significa CVE?}{}

\pergunta{Para que serve um identificador como CVE-2024-XXXX?}{}

\pergunta{Quem mantém o sistema CVE?}{}

\pergunta{Encontre uma CVE real e explique, em linguagem simples, qual era a vulnerabilidade.}{}

\pergunta{A vulnerabilidade encontrada recebeu alguma classificação de severidade?}{}

% ======================================================
\section{CVSS --- uma vulnerabilidade pode ser mais ou menos grave?}

\pergunta{Para que serve o CVSS?}{}

\pergunta{O que significa uma pontuação de 0 a 10?}{}

\pergunta{O que diferencia uma vulnerabilidade de baixa, média, alta ou crítica severidade?}{}

\end{document}
